\section{Readable four-vertex lower-bound certificates}
\label{app:n4-lower}

This appendix expands the lower-bound part of
Theorem~\ref{thm:n4-minimum}.  It deliberately stops at the ingredients needed
for the minima $7$ and $11$; it does not classify all equality families.

\subsection{The forced linear layer}

There is one degree-one support orbit: a single directed arc.  Let $P_0$ be
the empty pattern and $P_1$ a one-arc pattern.  Every coordinate of degree at
least two vanishes on both, while the linear orbit sum takes values zero and
one.  The separator mask is therefore a singleton, forcing the linear
coordinate in every separating family.

\subsection{The complete quadratic constraint graph}

Let $T_1,\ldots,T_r$ be canonical representatives of the quadratic support
orbits, where $r=5$ for $S_4$ and $r=7$ for $A_4$.  For $i\neq j$, compare the
two patterns $P_i=T_i$ and $P_j=T_j$.  Both have two arcs, so their unique
linear coordinate values agree.  Every cubic coordinate vanishes.  Among the
quadratic coordinates, $q_{[T_i]}$ and $q_{[T_j]}$ are the only coordinates
whose values distinguish the two support-orbit representatives.  Hence
\[
 \Sep(P_i,P_j)=\{q_{[T_i]},q_{[T_j]}\}.
\]
All $\binom r2$ pairs occur, producing the edge set of $K_r$.  A hitting set
must contain a vertex cover of that complete graph.  Since
$\tau(K_r)=r-1$, at least four of the five $S_4$ quadratics and six of the
seven $A_4$ quadratics are necessary.

\subsection{Coordinate-disjoint cubic constraints}

After the duplicate and superset reductions, the following explicit pattern
pairs have pure-cubic separator masks.  Within each action their coordinate
sets are pairwise disjoint.

\begin{center}
\small
\begin{tabularx}{\textwidth}{>{\raggedright\arraybackslash}p{0.12\textwidth}>{\raggedright\arraybackslash}p{0.18\textwidth}>{\raggedright\arraybackslash}X}
\toprule
Action & pattern pair & cubic separator coordinates\\
\midrule
$S_4$ & P111 / P155 & C001, C002, C006, C007\\
$S_4$ & P052 / P066 & C004, C005, C009, C011\\
\midrule
$A_4$ & P049 / P060 & C002, C005, C011, C016\\
$A_4$ & P112 / P115 & C006, C009, C014, C018\\
$A_4$ & P093 / P096 & C001, C008, C010, C019\\
$A_4$ & P106 / P153 & C000, C003, C004, C007, C012, C013, C017, C020\\
\bottomrule
\end{tabularx}
\end{center}

Here C$j$ abbreviates the coordinate ID \coord{S4-D3-C$j$} or
\coord{A4-D3-C$j$} in the corresponding row, with three decimal digits.  The
atlas stores the full pattern masks, values, pair IDs, and separator-mask
integers.  Pairwise disjointness makes the lower contribution two for $S_4$
and four for $A_4$.

Because the linear, quadratic, and cubic coordinate sets are disjoint, their
lower bounds add.  This proves $1+4+2=7$ and $1+6+4=11$ without assuming that
all lower-degree coordinates are retained.

\subsection{Uniform branch-DAG cross-check}

The complete reduced constraint antichains have $26$ and $58$ rows.  Starting
with budgets six and ten, a generic branch checker repeatedly applies forced
singleton and disjoint-packing rules and branches on an uncovered constraint.
The resulting proof DAGs have $11$ and $22$ nodes.  At every branch the listed
children cover every coordinate choice available in the selected constraint;
every leaf displays a packing larger than its remaining budget.  Independent
verification of these DAGs supplies a method-uniform check of the more
readable degree-layer argument above.
